% Modified from openai/math (Apache-2.0): new 11-strand proof of Proposition 3.1.
\section{A quantitative constraint from braid groups}\label{sec:braids}

Let $B_l$ denote the group on $s_1,\ldots,s_{l-1}$ with relations
\[
 s_i s_{i+1}s_i=s_{i+1}s_i s_{i+1},\qquad
 s_i s_j=s_j s_i\quad\text{if }|i-j|\ge2.
\]
Throughout this section the same notation may denote elements of another
group satisfying these relations.  No injectivity or cocompactness of
the resulting braid subgroup is needed for the length estimates.

\begin{proposition}\label{prop:braid-angle}
Suppose a group acting geometrically on a CAT(0) space contains elements
$s_1,\ldots,s_{10}$ satisfying the $B_{11}$ relations, with
$\ell(s_1^2)>0$.  Then
\begin{equation}\label{eq:braid-angle-bound}
 \left(\frac{\ell((s_1s_2)^3)}{3\ell(s_1^2)}\right)^2
 \ge\frac{7}{27}>\frac14,
\end{equation}
and in particular
\[
 \frac{\ell((s_1s_2)^3)}{3\ell(s_1^2)}
 =\frac{\ell((s_1s_2)^3)}{3\ell(s_2^2)}>\frac12.
\]
\end{proposition}

For $B_4$, Crisp and Paoluzzi compute translation components along an
axis of a central element under a minimal geometric action
\cite[Proposition~5]{CrispPaoluzzi2005}; minimality means that no proper
nonempty closed convex subset is invariant.

To prove Proposition~\ref{prop:braid-angle}, we compare central elements
of the initial braid subblocks using the centralizer characters of
Section~\ref{sec:lengths}.  For a general $B_l$ with $l\ge3$, define
\begin{equation}\label{eq:braid-centers-definition}
 L_1=1,\qquad L_j=s_{j-1}L_{j-1}s_{j-1},\qquad
 Z_r=L_1\cdots L_r,\qquad N_r=\binom r2.
\end{equation}
Here $2\le j\le l$ and $1\le r\le l$.

\begin{lemma}\label{lem:braid-centers}
The element $L_j$ commutes with $s_i$ for $i\le j-2$, and $Z_r$ commutes
with every $s_i$ for $i<r$.  In particular,
\[
 Z_2=s_1^2,\qquad Z_3=(s_1s_2)^3.
\]
The element $\delta=s_1\cdots s_{l-1}$ satisfies
\begin{equation}\label{eq:braid-shift}
 \delta s_i\delta^{-1}=s_{i+1}\qquad(1\le i\le l-2).
\end{equation}
In any geometric action of the ambient group,
\begin{equation}\label{eq:braid-central-pairing}
 \langle s_i^2,Z_r\rangle=\frac{\ell(Z_r)^2}{N_r}
 \qquad(2\le r\le l,\ 1\le i<r).
\end{equation}
\end{lemma}

\begin{proof}
Induct on $j$ for the assertion about $L_j$.  For $i\le j-3$, both
$s_{j-1}$ and $L_{j-1}$ commute with $s_i$.  For $i=j-2$, put
$a=s_{j-2}$, $b=s_{j-1}$, and $C=L_{j-2}$.  Then $bC=Cb$ and
$L_j=baCab$, so the braid relation gives
\[
 aL_j=abaCab=babCab=baCbab=baCaba=L_j a.
\]
This includes the case $j=3$, where $C=L_1=1$; the case $j=2$ is empty.
In particular, $L_j$ commutes with each earlier $L_i$, because the word
for $L_i$ uses only $s_1,\ldots,s_{i-1}$.

Next induct on $r$ for $Z_r$.  Both $Z_{r-1}$ and $L_r$ commute with
$s_i$ for $i\le r-2$.  The only remaining generator is $s=s_{r-1}$.
It commutes with $Z_{r-2}$.  With $L=L_{r-1}$ we have $L_r=sLs$,
and the already proved commutation of $L$ with $L_r$ says
$LsLs=sLsL$.  Therefore
\[
 s(LL_r)=sLsLs=LsLs^2=(LL_r)s,
\]
which proves the assertion.  The case $r=2$ is immediate.  Finally
\[
 Z_3=s_1^2s_2s_1^2s_2
     =s_1(s_1s_2s_1)s_1s_2
     =s_1(s_2s_1s_2)s_1s_2=(s_1s_2)^3.
\]

To verify \eqref{eq:braid-shift}, in $\delta s_i$ commute the last
$s_i$ left past $s_{i+2},\ldots,s_{l-1}$, replace
$s_i s_{i+1}s_i$ by $s_{i+1}s_i s_{i+1}$, and commute the new
$s_{i+1}$ left past $s_1,\ldots,s_{i-1}$.  The result is
$s_{i+1}\delta$.

For the length identity, apply the character
$h\mapsto\langle Z_r,h\rangle$ from
Lemma~\ref{lem:centralizer-character}.  All $s_i$, $i<r$, are in its
domain.  Applying it to each adjacent braid relation shows that their
values are equal; call the common value $\alpha$.  The positive word
for $L_j$ has $2(j-1)$ letters, so that for $Z_r$ it has $2N_r$ letters.
Thus
\[
 \ell(Z_r)^2=\langle Z_r,Z_r\rangle=2N_r\alpha,
 \qquad \langle s_i^2,Z_r\rangle=2\alpha,
\]
which proves \eqref{eq:braid-central-pairing}, including when
$\ell(Z_r)=0$.
\end{proof}


\begin{proof}[Proof of Proposition~\ref{prop:braid-angle}]
We work in $B_l$ with $l\ge4$ and specialize to $l=11$ only at the end.
Put $g=s_1^2$ and $D=\ell(g)>0$.  Since
$s_2=(s_1s_2)s_1(s_1s_2)^{-1}$, also $\ell(s_2^2)=D$.  For
$2\le r\le l$ put
\[
 Q_r=\frac{\langle g,Z_r\rangle}{D^2},\qquad\text{so that}\qquad
 \ell(Z_r)^2=N_rQ_rD^2
\]
by \eqref{eq:braid-central-pairing} with $i=1$.  Thus $Q_2=1$, and
$Q_3=\ell(Z_3)^2/(3D^2)=3R_3$, where
$R_3=\ell((s_1s_2)^3)^2/(9D^2)$ is the square of the ratio in
\eqref{eq:braid-angle-bound}.

\emph{A recursion for $Q_r$.}
All $s_j$ with $j\ge3$ commute with $g$.  Applying the character of $g$
(Lemma~\ref{lem:centralizer-character}) to the braid relation between
$s_j$ and $s_{j+1}$, $3\le j\le l-2$, shows that $\langle g,s_j\rangle$
is independent of $j\ge3$; write it as $c_0D^2/2$.  Put
$x=Q_3-1$; since $Z_3=gL_3$, we have $\langle g,L_3\rangle=xD^2$.
For $j>3$ the recursive definition gives
\[
 L_j=(s_{j-1}\cdots s_3)L_3(s_3\cdots s_{j-1}),
\]
and every displayed factor lies in $C_G(g)$ (for $L_3$ by
Lemma~\ref{lem:braid-centers}).  Keeping $L_3$ as one factor, the
character gives $\langle g,L_j\rangle=\bigl(x+(j-3)c_0\bigr)D^2$ for
$j\ge3$.  Summing over the factors $L_2=g,L_3,\ldots,L_r$ of $Z_r$,
\begin{equation}\label{eq:braid-recursion}
 Q_r=1+(r-2)x+\binom{r-2}{2}c_0\qquad(2\le r\le l).
\end{equation}
This calculation never expands the noncommuting $s_2$-letters inside
$L_3$.

\emph{Two positivity constraints.}
First, $\ell(Z_l)^2=N_lQ_lD^2\ge0$, so
\begin{equation}\label{eq:top-center}
 Q_l=1+(l-2)x+\binom{l-2}{2}c_0\ge0.
\end{equation}
Second, $Z_{l-1}$ and $Z_l$ commute, because every letter of the positive
word for $Z_{l-1}$ is some $s_i$ with $i\le l-2$.  Applying the
character of $Z_l$ to that word, as in the proof of
Lemma~\ref{lem:braid-centers}, each of its $2N_{l-1}$ letters
contributes $\ell(Z_l)^2/(2N_l)$, so
\[
 \langle Z_{l-1},Z_l\rangle=\frac{N_{l-1}}{N_l}\,\ell(Z_l)^2
 =N_{l-1}Q_lD^2.
\]
By Lemma~\ref{lem:commuting-form}, the element
$W=Z_{l-1}^{\,l}Z_l^{-(l-2)}$ has
\begin{align*}
 \frac{\ell(W)^2}{D^2}
 &=l^2N_{l-1}Q_{l-1}-2l(l-2)N_{l-1}Q_l+(l-2)^2N_lQ_l\\
 &=\frac{l(l-1)(l-2)}{2}\bigl(lQ_{l-1}-(l-2)Q_l\bigr),
\end{align*}
using $N_l=l(l-1)/2$ and $N_{l-1}=(l-1)(l-2)/2$.  Hence
$lQ_{l-1}-(l-2)Q_l\ge0$, which by \eqref{eq:braid-recursion} reads
\begin{equation}\label{eq:adjacent-centers}
 2+(l-4)x-2(l-3)c_0\ge0.
\end{equation}

\emph{Elimination of $c_0$.}
Add \eqref{eq:top-center} to $(l-2)/4$ times
\eqref{eq:adjacent-centers}.  The coefficients of $c_0$ are
$\binom{l-2}{2}$ and $-(l-2)(l-3)/2$, which cancel, leaving
\[
 \frac l2+\frac{l(l-2)}{4}\,x\ge0,\qquad\text{that is,}\qquad
 x\ge-\frac{2}{l-2}.
\]
Therefore
\begin{equation}\label{eq:general-braid-bound}
 R_3=\frac{1+x}{3}\ge\frac{l-4}{3(l-2)},
\end{equation}
which exceeds $1/4$ exactly when $l\ge11$.  For $l=11$ it gives
$R_3\ge 7/27>1/4$.  Finally $\ell(s_2^2)=\ell(s_1^2)=D$ gives the
equality of the two displayed ratios.
\end{proof}

\begin{remark}
The bound \eqref{eq:general-braid-bound} is the best obtainable from
\eqref{eq:braid-recursion}, \eqref{eq:top-center} and
\eqref{eq:adjacent-centers}: at $x=-2/(l-2)$ and
$c_0=2/((l-2)(l-3))$ one has $Q_{l-1}=Q_l=0$.  For $l=10$ it gives only
$R_3\ge1/4$, which is not enough for the strict sector inequality.
Numerically, adding the positive semidefiniteness of the full Gram
matrices of the commuting families built from disjoint and nested
braid subblocks does not improve \eqref{eq:general-braid-bound}.
\end{remark}
